\begin{table}[htbp]
\centering
\caption{Molecular top-1 of 30 over the leave-one-agent-out calibration envelope,
scanned continuously on the 0.002 grid. The eleven held-out refits of
Table~\ref{tab:si-loao} span $s_\mathrm{low}$ = 1.006--1.024 and
$s_\mathrm{high}$ = 0.980--0.984 but land on only four distinct values of
$s_\mathrm{low}$, leaving the interior of their own range unsampled. The adopted
calibration is in bold. Produced by \texttt{code/v3\_scalescan.py}.}
\label{tab:si-scalescan}
\small
\begin{tabular}{lccc}
\toprule
 & \multicolumn{3}{c}{$s_\mathrm{high}$} \\
\cmidrule(lr){2-4}
$s_\mathrm{low}$ & 0.980 & 0.982 & 0.984 \\
\midrule
1.006 & 13 & 13 & 13 \\
1.008 & 13 & 13 & 13 \\
1.010 & 13 & 13 & 13 \\
1.012 & 13 & 13 & 13 \\
\textbf{1.014} & \textbf{13} & 13 & 13 \\
1.016 & 16 & 16 & 16 \\
1.018 & 16 & 16 & 16 \\
1.020 & 12 & 12 & 13 \\
1.022 & 12 & 12 & 12 \\
1.024 & 12 & 12 & 12 \\
\bottomrule
\end{tabular}

\begin{flushleft}\footnotesize The molecular figure runs from 12 to 16 of 30 across
the envelope and the adopted calibration returns 13. $s_\mathrm{high}$ moves it at
one point only; the variation is carried entirely by $s_\mathrm{low}$. The maximum at
$s_\mathrm{low}$ = 1.016 and 1.018 is reached because all three HN1 acquisitions move
from HN3 to HN1; those values lie 0.2 and 0.4\,\% from the adopted 1.014 and inside
the held-out range, so the adopted figure is not the most favorable one available
within the calibration's own uncertainty. We report it because it is the value the
non-circular objective selects, and we state the envelope rather than a robustness
claim the eleven-point check cannot support.\end{flushleft}
\end{table}
